\documentclass[11pt]{article}
\usepackage{icmlko}
\begin{document}

\title{한 수로는 안 나온다 --- 축의 몫이 도메인마다 $+98\%$ 에서
$-67\%$ 다}
\author{월드모델 실험실}
\date{노트 481 · 2026년 8월 3일}
\maketitle

\begin{abstract}
노트 480 이 스스로 적어 둔 흠을 고쳤다. 그 노트의 대조는 \emph{전체
기준} 덮음 $95\%$ 축 셋을 썼는데, 그 셋이 \textbf{도메인 안에서는 고르게
없었다} --- 모바일에서 대조 격차가 전체보다 컸던 것이 증거였다.

도메인마다 \emph{그 도메인의} 유보에서 덮음 $\ge 95\%$ 인 축만 골라
대조를 다시 만들었다(대조 축의 두 절반 덮음이 $99\sim100\%$ 임을 확인).

\begin{center}
\begin{tabular}{lrr}
\hline
가중평균 격차 & 전체 축 & 맞춘 대조 \\
\hline
 & $+0.1028$ & $\mathbf{+0.0415}$ \\
\hline
축의 몫 & \multicolumn{2}{c}{$\mathbf{60\%}$ \quad(노트 480 은 $21\%$)} \\
\hline
\end{tabular}
\end{center}

\textbf{샌 대조가 축을 과소평가하고 있었다.} $21\%$ 가 아니라 $60\%$ 다.

\textbf{그런데 등록은 반증이다} --- 모바일의 이상이 안 사라졌다(대조
$+0.2497$ 대 전체 $+0.2244$). 그리고 도메인별로 열어 보니 \textbf{$60\%$
라는 수 자체가 못 쓸 요약}이었다.

\begin{center}
\begin{tabular}{lrrr}
\hline
도메인 & 전체 격차 & 맞춘 대조 & 축의 몫 \\
\hline
게임 & $+0.1180$ & $+0.0018$ & $\mathbf{+98\%}$ \\
만화 & $+0.2640$ & $+0.0962$ & $+64\%$ \\
세계애니 & $+0.1865$ & $+0.1352$ & $+28\%$ \\
모바일 & $+0.2244$ & $+0.2497$ & $-11\%$ \\
애니 & $-0.1004$ & $-0.1678$ & $\mathbf{-67\%}$ \\
\hline
\end{tabular}
\end{center}
\end{abstract}

\begin{figure}[t]
\centering
\includegraphics[width=\textwidth]{figs/f1.pdf}
\caption{\textbf{가} --- 도메인마다 맞춘 대조. \textbf{나} --- 축의 몫이
$+98\%$ 에서 $-67\%$ 까지 퍼져 있다. 가중평균($60\%$)도 노트 480 의
$21\%$ 도 어느 도메인도 안 닮았다.}
\caption*{}
\end{figure}

\section{게임과 애니가 정반대다}

\textbf{게임} --- 전체 격차 $+0.1180$ 이 맞춘 대조에서 $+0.0018$ 로
\emph{거의 사라진다}. 게임에서 축 많은 행이 잘 맞는 것은 \textbf{순전히
축 때문}이다. 게임은 자기만 갖는 축이 넷 있고(가격 · 연령등급 · 램 ·
카테고리 수) 그것들이 있고 없고가 그대로 성능이 된다.

\textbf{애니} --- 격차가 $-0.1004$ 로 애초에 \emph{반대}이고 맞춘
대조에서 $-0.1678$ 로 \textbf{더 벌어진다}. 애니에서는 축이 적은 행이
더 잘 맞고, 축을 줄이면 그 우위가 커진다.

\textbf{한 수로 요약하면 이 둘이 상쇄된다.} 가중평균 $60\%$ 는 게임의
$98\%$ 도 애니의 $-67\%$ 도 안 닮았다.

\section*{조항 --- 폭을 안 적으면 평균을 안 적는다}

\begin{quote}
\textbf{조항} --- 도메인별로 잰 것을 하나로 요약할 때는 \textbf{폭을 같이
적는다.} 부호가 갈리면(여기서는 $+98\%$ 와 $-67\%$) \textbf{평균을 결론
문장에 쓰지 않는다.} 노트 480 의 ``$21\%$'' 와 이 노트의 ``$60\%$'' 는
둘 다 그 규칙을 어긴 수다 --- 앞엣것은 대조가 새서, 뒤엣것은 폭이 커서.
\end{quote}

\section{태거 갈래는 어떻게 되나}

노트 480 이 태거의 기대값을 내렸는데, 이 노트는 \textbf{그것을 도메인별로
가른다.}

\begin{itemize}
\item \textbf{게임처럼 몫이 $98\%$ 인 곳} --- 축을 더 붙이면 그대로
번다. 게임의 유보 격차 $+0.1180$ 이 거의 전부 축의 값이다.
\item \textbf{애니처럼 음수인 곳} --- 축을 더 붙여도 이 격차는 안 준다.
다른 이야기가 있다.
\end{itemize}

\emph{``태거를 늘려라''는 도메인을 골라서 할 말이다.}

\section*{남은 것}

\begin{itemize}
\item \textbf{애니가 왜 반대인가} --- 두 노트째 미해결. 축이 적은 행이
더 잘 맞고, 축을 줄이면 더 잘 맞는다.
\item \textbf{모바일의 이상도 대조 탓이 아니었다} --- 덮음을 맞췄는데도
남는다. 남은 설명은 \emph{대조 팔이 다시 학습한다}는 것이다(같은 모형이
아니라 다른 모형이다). 그 설계를 고치려면 예보 때만 가려야 하는데,
그러면 모델이 없는 축을 기대하는 문제가 생긴다.
\item \textbf{게임의 전용 축 넷} --- 몫 $98\%$ 의 정체가 그것인지 직접
확인 안 했다.
\end{itemize}

\end{document}
